\subsection{Exact classical specialization of the ordinary modules}
\label{sec:selmer-classical-control}

Fix the original definite unitary datum, type lattice, auxiliary set,
uniformizers, and unnormalized Hecke operators.  Write
\(A=\mathcal O[[X_1,\ldots,X_D]]\), \(D=n[F:\mathbb Q]\), for its
weight ring, and \(M(K,A,\theta)\) for the ordinary function module.
The preceding integral ordinary construction supplies finite freeness,
exact reduction, and the ordinary projector.  These facts give the
following actual specialization theorem.

\begin{proposition}[Coefficient change and classical control]
\label{prop:selmer-classical-control}
For every continuous local map \(f:A\to B\) of \(\mathcal O\)-algebras,
with \(B\) complete Noetherian local and with finite residue field,
pointwise coefficient change gives
\[
 M(K,A,\theta)\otimes_A B\simeq M(K,B,f\theta).
\]
This respects the original Hecke operators, auxiliary translations,
diamond actions, traces, and prescribed residual Hecke summand.
If \(B=\mathcal O_E\) for a finite coefficient extension and \(f\theta\)
has finite image, these modules specialize faithfully to the classical
weight-zero ordinary spaces at every level \(c\geq c_0(f)\), where
\(c_0(f)\) is independent of \(K\).  The specialized basis is a basis
of that space.
\end{proposition}

\begin{proof}
Given \(k\), choose \(N\) with
\(f(\mathfrak m_A^N)\subseteq\mathfrak m_B^k\), and choose \(c\) such
that \(\theta(T_c)=1\) modulo \(\mathfrak m_A^N\).  At the same
original double-coset representatives, the finite function spaces are
free coordinate modules, and
\[
 E(K,A/\mathfrak m_A^N,c)\otimes_A B/\mathfrak m_B^k
       =E(K,B/\mathfrak m_B^k,c).
\]
The integral sections and trivial stabilizers from the ordinary
construction justify this equality; no flatness assumption or
commutation of arbitrary invariants with base change is needed.
The matrix of the original unnormalized operator
\[
 U_z=\sum_{u\in N/(zNz^{-1})}[uz],
 \qquad z=\operatorname{diag}(\varpi^{n-1},\ldots,\varpi,1),
\]
and of its product over the places above \(\ell\), changes coefficients
by \(f\).

The finite Fitting decomposition into an invertible and a nilpotent
part is split.  Tensoring preserves the direct sum, invertibility,
and nilpotence, hence gives the Fitting decomposition after coefficient
change.  Exact reduction of the ordinary module therefore yields
\[
 M(K,A,\theta)\otimes_A B/\mathfrak m_B^k
       \simeq M(K,B,f\theta)/\mathfrak m_B^k.
\]
The maps are pointwise and compatible.  Taking inverse limits proves
the assertion, since the original module is finite free.
All operator identities follow from their finite-sum formulas.

For residual localization, the finite commutative Hecke algebra has
an idempotent factor on which every prescribed \(H_i-a_i\) is
topologically nilpotent, while on the complementary factor these
elements generate the unit ideal.  A finite set of operators suffices
by Noetherianity.  Both conditions persist under local coefficient
change, so the original residual idempotent is preserved.
The Hecke family here remains the original spherical family outside
the fixed \(S\cup Q\), together with the original positive diagonal
operators above \(\ell\).

Now suppose \(f\theta(T_{c_0})=1\) exactly.  For every \(k\) and
\(c\geq c_0\), level lowering places the ordinary reduction inside
\(E(K,\mathcal O_E/\mathfrak m_E^k,c)\).  Thus an integral ordinary
vector is exactly \(\overline N(c)\)-invariant: every invariance
difference belongs to all powers of \(\mathfrak m_E\).  It lies in
the finite classical lattice \(E(K,\mathcal O_E,c)\), with the
original type, tame level, \(N\)-invariance, and inverse torus
character.

On this lattice the integral ordinary decomposition has \(U\) and
\(U^{-1}\) integral on its first summand and \(U\) nilpotent modulo
\(\mathfrak m_E\) on the second.  Over
\(\overline{\mathbb Q}_\ell\), these are exactly the generalized
eigenspaces with eigenvalues of absolute value \(1\) and \(<1\).
At weight zero the algebraic representation is trivial,
\(\operatorname{lowW}=0\), and all source
\(\operatorname{algPow}\) and \(\operatorname{lowVal}\) factors are
\(1\).  Thus neither \(U\) nor any Hecke normalization changes.

Apply \(\iota f\) pointwise, for the fixed
\(\iota:\overline{\mathbb Q}_\ell\simeq\mathbb C\).
Coefficient change proves independence of the specialized ordinary
basis and that it spans the entire unit-slope space.  The classical
space is finite-dimensional, its Hecke operators preserve the
integral lattice and have integral eigenvalues, and all semilinear,
Hecke, and auxiliary-action identities hold term by term.  This is
the complete weight-zero specialization interface.
\end{proof}

\begin{proposition}[Enough good points at weight zero]
\label{prop:selmer-good-zero}
The original weight ring has a family of finite-order classical
specializations with pairwise distinct local torus characters which
detects every nonzero element of \(A\).  They detect every nonzero
vector in the original ordinary modules, their prescribed residual
summands, and their smooth auxiliary union.
\end{proposition}

\begin{proof}
Local reciprocity identifies the free pro-\(\ell\) inertial quotient
with \(\Gamma_v\simeq\mathbb Z_\ell^{d_v}\), where
\(d_v=[F_v:\mathbb Q_\ell]>0\); see
\cite{MilneClassicalControl}.  Its rank follows directly from
logarithm and exponential on sufficiently deep principal units.
Keep the original fixed characters on the finite torsion subgroup.
Choose a finite-order lift \(\psi_{0,v,i}\) of each residual local
character with precisely that torsion restriction: combine its
Teichmuller lift with the prescribed finite torsion character using
a splitting of the local abelian group.

Write the original cyclotomic character as
\(\varepsilon=\omega_\varepsilon\langle\varepsilon\rangle\).
Here \(\omega_\varepsilon\) has order dividing \(\ell-1\), while
\(\langle\varepsilon\rangle\) takes values in
\(1+\ell\mathbb Z_\ell\).  Since \(\ell\) is odd, the latter group
is torsion-free.  For arbitrary finite-order
\(\chi_{v,i}:\Gamma_v\to\mu_{\ell^\infty}\), set
\[
 \kappa_{v,i}^x
    =\psi_{0,v,i}\langle\varepsilon\rangle^{-i}\chi_{v,i}
       \quad\text{on inertia}.
\]
Extend \(\chi_{v,i}\) to the local decomposition group with value
\(1\) on a chosen uniformizer.  The residual and prescribed torsion
restrictions are unchanged, so the original universal weight
property gives \(x:A\to\mathcal O_E\) for a finite extension \(E\).
This selects no unramified ordinary eigenvalue.

The exact source identity
\[
 \kappa_{v,i}(g)\varepsilon(g)^i\theta(e_{v,i}(y))=1,
 \qquad\operatorname{Art}(g)=y,
\]
now gives
\[
 x\theta(e_{v,i}(y))
   =\bigl(\psi_{0,v,i}(g)\omega_\varepsilon(g)^i
                       \chi_{v,i}(g)\bigr)^{-1}.
\]
It has finite order.  This principal cyclotomic twist retains the
source's weight shift; unshifted torsion weights would not suffice.

Take the original coordinates
\[
 1+X_{v,i,j}
   =\psi_{0,v,i}(\gamma_{v,j})^{-1}
                 \kappa_{v,i}(\gamma_{v,j}),
\]
with \(\gamma_{v,j}\) projecting to a basis of \(\Gamma_v\).
At the points just constructed,
\[
 1+x(X_{v,i,j})=u_{v,i,j}\zeta_{v,i,j},\qquad
 u_{v,i,j}=\langle\varepsilon(\gamma_{v,j})\rangle^{-i}
     \in1+\mathfrak m_{\mathcal O},\quad
 \zeta_{v,i,j}\in\mu_{\ell^\infty}.
\]
The substitution \(X=u(1+T)-1\) is an automorphism of the complete
power-series ring.

For a nonzero \(h(T)\in\mathcal O_E[[T]]\), remove the largest common
uniformizer power from its coefficients and let \(r\) be the first
unit coefficient.  If \(0<v_E(t)\) and \(rv_E(t)<1\), the \(r\)-th
term is the unique term of least valuation, so \(h(t)\ne0\).
Primitive roots of increasing \(\ell\)-power order have
\[
 v_E(\zeta_{\ell^a}-1)
   =\frac{v_E(\ell)}{\ell^{a-1}(\ell-1)}\longrightarrow0.
\]
For \(r=0\) the same conclusion needs no upper bound on the valuation.
In several variables, keep one nonzero coefficient in the last
variable, specialize its other variables inductively, and then
apply this one-variable argument.  All partial specializations lie
in finite coefficient extensions.

This induction also permits finitely many inequalities
\(\alpha_s\zeta_s\ne\alpha_t\zeta_t\): each earlier choice forbids
at most one value of the later root.  At one basis element for each
\(\Gamma_v\), take
\(\alpha_{v,i}=\psi_{0,v,i}(\gamma_{v,1})
\omega_\varepsilon(\gamma_{v,1})^i\).
The displayed local character formula then makes the characters
pairwise distinct.  The source admits weight zero, so these points
are good in its precise sense.

Finally choose a nonzero coordinate of any nonzero ordinary vector
in an \(A\)-basis.  A good point detects it; the first proposition
and independence of the specialized basis detect the vector.
A vector in the smooth union lies at one compact auxiliary level.
\end{proof}

\paragraph{Exact boundary.}
The source's actual Hecke-ideal construction chooses weight zero,
so these propositions supply its specialization and dense-weight
inputs \cite{SelmerClassicalControlSource}.  They do not yet prove
that classical Hecke eigencharacters separate the Hecke algebra,
construct its continuous ordinary Galois determinant, or supply
the full Galois-compatible classical-point interface.  In
particular W04, its patched support and uniform integral Selmer
bounds remain open.  The field degrees, \(D=n[F:\mathbb Q]\), and
specialization conductors may grow also in the ordinary RH route.
